% concatenated from random-Weierstrass-function-V6.tex
\documentclass[12pt, reqno]{amsart}
\usepackage{amsmath,amssymb,amsthm}
\usepackage{mathrsfs}
\usepackage{amsfonts}
\usepackage{hyperref}
\usepackage{enumitem}
\usepackage{graphicx}
\usepackage{mathrsfs}
\usepackage{upgreek}
\usepackage{mathtools}
\usepackage{tikz}       % 核心：绘图必须加载
\usepackage{amsmath}    % 数学公式兼容
\setlength{\textwidth}{15.0cm}
\setlength{\textheight}{22.0cm}
\hoffset=-1cm
\pagestyle {plain}
 \topmargin=0cm
\errorcontextlines=0

% 定理环境
\theoremstyle{plain}
\newtheorem{theorem}{Theorem}[section]
\newtheorem{proposition}{Proposition}[section]
\newtheorem{lemma}{Lemma}[section]
\newtheorem{corollary}{Corollary}[section]

\theoremstyle{definition}
\newtheorem{definition}{Definition}[section]
\newtheorem{example}{Example}[section]

\theoremstyle{remark}
\newtheorem{remark}{Remark}[section]

% 简写
\newcommand{\R}{\mathbb{R}}
\newcommand{\N}{\mathbb{N}}
\newcommand{\Z}{\mathbb{Z}}
\newcommand{\C}{\mathbb{C}}
\newcommand{\PP}{\mathbb{P}}
\newcommand{\M}{\bar{M}}
\newcommand{\eps}{\varepsilon}
\newcommand{\del}{\partial}
\newcommand{\la}{\langle}
\newcommand{\ra}{\rangle}
\newcommand{\td}{\tilde}
\newcommand{\htd}{\hat}
\newcommand{\ol}{\overline}

\begin{document}
	
\title{On the Hausdorff dimension of graph of random vector-valued Weierstrass function}

\author{Jun Jason Luo}  \address{College of Mathematics and Statistics,  Chongqing University, Chongqing, 401331, P.R. China} \email{jun.luo@cqu.edu.cn}


\author{Zi-Rui Zhang}  \address{College of Mathematics and Statistics,  Chongqing University, Chongqing, 401331, P.R. China} \email{202306021017@stu.cqu.edu.cn}

%\author{Peng-Fei Zhang} \address{Department of Mathematics, Southwest Jiaotong University, Chengdu, 611756, P.R. China}
%\email{pfzhang@link.cuhk.edu.hk}

\begin{abstract}
		Let $\Theta=\{\theta_n\}, \Lambda=\{\lambda_n\}$ be two sequences of independent and identically distributed uniform random variables on $[0,1]$. The  random vector-valued Weierstrass function is given by
		$$
		f_{\Theta,\Lambda}(x)=
		\left(
		\sum_{n=0}^{\infty} a^n\cos\bigl(2\pi (b^n x+\theta_n)\bigr),\ 
		\sum_{n=0}^{\infty} a^n\sin\bigl(2\pi (b^n x+\lambda_n)\bigr)
		\right), \;  x\in[0,1],
		$$ where $0<a<1<b,\ ab> 1$. The Hausdorff dimension of the graph of this function is proved to be
	$$\dim_H G(f_{\Theta,\Lambda}) = \min\left\{-\frac{\log b}{\log a}, \, 3 +2\frac{\log a}{\log b}\right\} \quad \text{a.s.}$$
\end{abstract}

	\date{\today}
	
	\keywords{Hausdorff dimension; Box dimension; Weierstrass function}
	
	\subjclass[2020]{Primary 28A78; Secondary 37C45}
	
	\thanks{The research was supported by the Natural Science Foundation of Chongqing (No.CSTB2023NSCQ-MSX0553)}
	
	
	\maketitle
	
	\section{Introduction}
	
In 1872, Weierstrass \cite{Weierstrass1872} constructed the pathological function in series form
\begin{equation}\label{w-cos function}
	W(x)=\sum_{n=0}^\infty a^n\cos(2\pi b^n x),\quad x\in [0,1],
\end{equation}
where $0<a<1<b,\ ab\ge 1$, and gave the first rigorous proof that $W$ is continuous everywhere but nowhere differentiable for some $a$ and $b$. This overturned the conventional intuition in classical analysis that ``continuity implies almost everywhere differentiability''. Hardy \cite{Hardy1916} refined the proof of non-differentiability and laid a solid theoretical foundation. In 1977, Mandelbrot \cite{Mandelbrot1977} proposed the famous conjecture:
\[
\dim_H(G(W))=\dim_B(G(W))=2+\frac{\log a}{\log b}
\] where $\dim_H, \dim_B$ are Hausdorff and box dimensions and $G(W)$ denotes the graph of $W$ over $[0,1]$. This conjecture has become an iconic open problem in fractal geometry.

Kaplan, Mallet-Paret, and Yorke \cite{Kaplan1984} employed methods from dynamical systems, viewing the graph of the Weierstrass function as the attractor of an expanding dynamical system, and established the box dimension formula mentioned above. Mauldin and Williams \cite{Mauldin1986} proved that $\dim_H(G(W))$ is bounded below by $2+{\log a}/{\log b}-O(1/{\log b})$ for sufficiently large $b$. Later, Przytycki and Urbański \cite{Przytycki1989} showed that $\dim_H(G(W))>1$. Ledrappier \cite{Ledrappier1992} used Pesin theory to obtain almost-sure dimension results for the sawtooth-type Weierstrass function.  Solomyak \cite{Solomyak1995} provided crucial support for the conjecture using transversality arguments and absolute continuity of Bernoulli convolutions. 

In related work, Hunt \cite{Hunt1998} constructed the randomized Weierstrass function by introducing an independent and identically distributed (i.i.d.) random phase sequence \(\Theta=\{\theta_n\}\):
\[
W_{\Theta}(x) = \sum_{n=0}^{\infty} a^n \cos\left(2\pi \left(b^n x + \theta_n\right)\right).
\]
He proved that with probability one, the Hausdorff dimension of the graph of $W_{\Theta}$ equals the conjectured value, i.e.,
\[
\dim_H(G(W_{\Theta}))=2+\frac{\log a}{\log b} \qquad \text{a.s.}
\]
An analogous result was established in \cite{Rom2014} by Romanowska through randomizing the parameter $a$. Heurteaux \cite{Heurteaux2003} studied the oscillation of  the following more general Weierstrass-type function with random phases
\[
f_{a,b}^{\phi,\Theta}(x) = \sum_{n=0}^\infty a^n \phi(b^n x + \theta_n)
\]
where $\phi$ is a Lipschitz function with period $1$. He showed that the function satisfies a H\"older condition almost surely, and both the packing dimension and box dimension of its graph equal $2 + {\log a}/{\log b}$ almost surely.  Baránski \cite{Baranski2012} extended the result to the   deterministic Weierstrass-type function $f^\phi(x) = \sum_{n=0}^\infty a_n \phi(b_n x +\theta_n)$ and determined  both the Hausdorff and box dimensions of its graph, with the dimension depending on the rapidly growing frequencies satisfying \(a_{n+1}/a_n\to 0\) and \(b_{n+1}/b_n\to \infty\).


A substantial advance toward Mandelbrot’s conjecture was achieved in \cite{Baranski2014}, where Baránski, Bárány and Romanowska verified the conjecture for all integers \(b>1\) over a large parameter class of $a$. Soon later, Shen \cite{Shen2018} fully settled the conjecture via tools from ergodic theory and dynamical systems, assuming integer \(b>1\) and \(1/b<a<1\). Ren and Shen \cite{Ren2021} further extended the preceding result. For every integer \(b>1\) with \(1/b<a<1\) and any real-analytic periodic function \(\phi:\mathbb{R}\to\mathbb{R}\), by defining 
\begin{equation}\label{w-function 1}
	f_{a,b}^\phi(x) = \sum_{n=0}^\infty a^n \phi(b^n x), 
\end{equation} they proved that the function \(f_{a,b}^\phi\) falls into exactly two categories: it is either real-analytic, whose graph has Hausdorff dimension $1$, or fractal, with its graph possessing Hausdorff dimension \(2+{\log a}/{\log b}\).

Meanwhile, partial progress has been achieved concerning the dimensional properties of graphs of vector-valued Weierstrass functions. It is known that when $b>1$ is an integer, the classical Weierstrass function \eqref{w-cos function} is the real part of the lacunary complex power series
\[
w(z) = \sum_{n=0}^{\infty} a^n z^{b^n}, \quad z\in \mathbb{C},\ |z| \leq 1
\]
on the unit circle. Salem and Zygmund \cite{Salem1945} and  Kahane, Weiss, and Weiss \cite{Kahane1963} showed that when $a$ is sufficiently close to $1$, the image of $w$ on the unit circle becomes a Peano curve. Baránski \cite{Baranski2002} further proved that the box dimension of the vector-valued Weierstrass function
\[
f_{a, b}(x)= \left( \sum_{n=0}^{\infty} a^n \cos(2\pi b^n x), \,  \sum_{n=0}^{\infty} a^n \sin(2\pi b^n x) \right)
\]
satisfies
\[
\dim_B(G(f_{a,b})) = 3 + 2\frac{\log a }{\log b},
\]
where the graph $G(f_{a,b})$ is a subset of $\mathbb{R}^3$ (see Figure \ref{fig1}). Ren \cite{Ren2023} extended Baránski’s result by replacing trigonometric functions with a vector-valued Lipschitz function $\phi:\mathbb{R}\to\mathbb{R}^2$ of period $1$. He proved that if $\mathbb{R}^2\setminus\phi(\mathbb{R})$ is disconnected, then for any integer $b>1$ and for $a$ sufficiently close to $1$, the box dimension of the graph of the vector-valued Weierstrass-type function of the form \eqref{w-function 1} is also given by $3 + 2{\log a }/{\log b}.$
More recently, Ren and Shen \cite{RenShen2024} improved their one-dimensional result into high-dimensional vector-valued Weierstrass functions.
	
		\begin{figure}[h] \label{fig1}
		\includegraphics[width=6.5cm]{weierstrass.pdf}
		\caption{The graph  $G(f_{a,b})$ where $a=0.6, b=2$.} 
	\end{figure}
	
Inspired by the aforementioned results, this paper is devoted to a further investigation of randomized vector-valued Weierstrass functions. Specifically, we focus on verifying the validity of the Mandelbrot's conjecture in high-dimensional random contexts, as well as establishing the consistency between the Hausdorff dimension and the box dimension of the graph of such functions.

Let \( H = [0,1]^\infty \) be the infinite-dimensional product space, and let
\( \Theta = \{\theta_n\} \in H \) and \( \Lambda = \{\lambda_n\} \in H \)
be two sequences of random phases, which are mutually independent and uniformly distributed over \([0,1]\). Based on these random phases, we introduce the two-dimensional randomized vector-valued Weierstrass function
\( f_{\Theta,\Lambda}: [0,1] \to \mathbb{R}^2 \) as follows:
\begin{equation}\label{random-vector-valued function}
f_{\Theta,\Lambda}(x) = \left(
\sum_{n=0}^{\infty} a^n \cos\bigl(2\pi (b^n x + \theta_n)\bigr),\,  \sum_{n=0}^{\infty} a^n \sin\bigl(2\pi (b^n x + \lambda_n)\bigr)
\right),
\end{equation}
where $ 0<a<1<b,\ ab> 1$. This function serves as a complex extension and randomized generalization of the classical one-dimensional Weierstrass function. Its graph is a typical self-affine fractal set in  $\mathbb{R}^3$, which combines randomness, self-affinity, and nowhere differentiability, making it an ideal object for testing fractal dimension theory.

For the above random fractal function, the main dimension result is stated as follows:

\begin{theorem}\label{main thm}
	Let \( f_{\Theta,\Lambda}: [0,1] \to \mathbb{R}^2 \) be the random vector-valued Weierstrass function defined in \eqref{random-vector-valued function}. The Hausdorff dimension of the graph satisfies
	\[
	\dim_H G(f_{\Theta,\Lambda}) =\dim_B G(f_{\Theta,\Lambda}) = \min\left\{-\frac{\log b}{\log a}, \, 3 +2\frac{\log a}{\log b}\right\} \quad \text{a.s.}
	\]
\end{theorem}

The upper bound for the Hausdorff dimension of the graph is known to equal the box dimension, which follows directly from the H\"older continuity of the function \(f_{\Theta,\Lambda}\). The main difficulty lies in establishing the corresponding lower bound. Let \(\beta = -{\log a}/{\log b}\). We split the proof into two cases: For \(\beta \in (0, \frac{1}{2})\), we adopt the potential-theoretic method introduced by Hunt \cite{Hunt1998}; For \(\beta \in [\frac{1}{2}, 1)\), we first compute the Hausdorff dimension of the image of \(f_{\Theta,\Lambda}\), and then use a projection argument to obtain the required lower bound for the graph. For a more detailed discussion on the Hausdorff dimension of images and graphs of random complex series, we refer the reader to the recent work of Lai, Lau, and Zhang \cite{LaiLauZhang2026}.
	
\section{Proof of Theorem \ref{main thm}}
For convenience, we denote by $\beta = - {\log a }/{\log b}.$ Obviously,  $\beta \in (0,1)$ and $a = b^{-\beta}$. The random vector-valued Weierstrass  function in \eqref{random-vector-valued function} can be rewritten as: for $x\in [0,1]$,
	\begin{equation*}
		f_{\Theta,\Lambda}(x) = \left( \sum_{n=0}^{\infty} b^{-\beta n} \cos 2\pi (b^n x + \theta_n),\,  \sum_{n=0}^{\infty} b^{-\beta n} \sin 2\pi (b^n x + \lambda_n) \right).
	\end{equation*}
More precisely, let
\begin{equation}\label{eq-phi-n}
	\varphi_n:=\varphi_n(x,y)=2b^{-\beta n}\sin(\pi b^n(x-y)). 
\end{equation} The sum-to-product formulas in trigonometry imply that
\begin{equation}\label{eq- vector-valued funct.}
	f_{\Theta,\Lambda}(x) -f_{\Theta,\Lambda}(y)=(a_1(x,y), a_2(x,y))
\end{equation}
 where 	
	\begin{eqnarray*}
	a_1(x,y) &=& \sum_{n=0}^{\infty} -\varphi_n \sin\left(\pi\left(b^n x + b^n y +2\theta_n\right)\right), \\
	a_2(x,y) &=& \sum_{n=0}^{\infty} \varphi_n \cos\left(\pi\left(b^n x + b^n y +2\lambda_n\right)\right).
   \end{eqnarray*}
	
Define the graph function of the function $f_{\Theta,\Lambda}$ by $$G(f_{\Theta,\Lambda}, x)= (x,  f_{\Theta,\Lambda}(x)): [0,1]\to \R^3.$$  Then the graph of  $f_{\Theta,\Lambda}$ over $[0,1]$ is given by $$G(f_{\Theta,\Lambda})=G(f_{\Theta,\Lambda}, [0,1]).$$ 

The following two results are straightforward, we provide proofs here for completeness.
	
	\begin{lemma}\label{lem-Holder}
	Both $f_{\Theta,\Lambda}(x)$ and $G(f_{\Theta,\Lambda},x)$ are $\beta$-H\"older continuous on $[0,1]$.
	\end{lemma}
	
	\begin{proof}
		For any \( x, y \in [0,1], x\ne  y\),
		\begin{eqnarray*}
				|f_{\Theta,\Lambda}(x) -f_{\Theta,\Lambda}(y)| &=& \sqrt{(a_1(x,y))^2+(a_2(x,y))^2} \\				 
				&\leq& 4 \sum_{n=0}^{\infty}  b^{-\beta n} \min\left\{1, \pi b^n |x - y|\right\} \\
				&\leq& 4  \pi \sum_{n=0}^{m-1} b^{(1-\beta) n} |x - y| + 4 \sum_{n=m}^{\infty} b^{-\beta n} \\
				&=& 4\pi \frac{b^{(1-\beta) m} - 1}{b^{1-\beta} - 1} |x - y| + 4 \frac{b^{-\beta m}}{1 - b^{-\beta}}
		\end{eqnarray*}
		for all \( m \ge 1 \).  Let \( m \) be the positive integer satisfying \( b^{-m} < |x - y| \leq b^{-m+1} \). It follows that
		\begin{eqnarray}\label{equ- constant}
			|f_{\Theta,\Lambda}(x) -f_{\Theta,\Lambda}(y)| &\leq& 4\pi \frac{(b / |x - y|)^{1-\beta}}{b^{1-\beta} - 1} |x - y| + 4\frac{|x - y|^{\beta}}{1 - b^{-\beta}} \nonumber \\
				&\leq& \left( \frac{4\pi b^{1-\beta}}{b^{1-\beta} - 1} + \frac{4}{1 - b^{-\beta}} \right) |x - y|^{\beta} \nonumber \\
				&:=& C |x - y|^{\beta}. 
		\end{eqnarray}		
		Moreover, $$|G(f_{\Theta,\Lambda},x)-G(f_{\Theta,\Lambda},y)|\le|x-y|+|f_{\Theta,\Lambda}(x) -f_{\Theta,\Lambda}(y)|\le (C+1)|x-y|^\beta.$$		
	\end{proof}
	

	
	\begin{lemma}\label{lem-lowerbound}
		$\dim_H G(f_{\Theta,\Lambda})\le \min\left\{\frac{1}{\beta}, \, 3-2\beta\right\}$ a.s.
	\end{lemma}
	
	\begin{proof}
	The $\beta$-H\"older continuity in	Lemma \ref{lem-Holder} yields that 
	$$\dim_H G(f_{\Theta,\Lambda})\le \frac{1}{\beta} \dim_H([0,1])=\frac{1}{\beta}.$$
		
	On the other hand,	take an integer \( m>1 \) and divide \( [0,1] \) into \( m \) subintervals of equal length. Let $J$ be one of the subintervals, H\"older continuity implies that
		\[
		\sup_{x,y\in J}|f_{\Theta,\Lambda}(x)-f_{\Theta,\Lambda}(y)| \le C m^{-\beta}
		\] where $C$ is the coefficient in \eqref{equ- constant}. 	The graph of the function on $J$ can be covered by a cuboid of size
		\( m^{-1}\times C m^{-\beta}\times C m^{-\beta} \), and hence by at most $(C m^{1-\beta} + 1)^2$ cubes of side length \( 1/m \). Thus the whole graph  $G(f_{\Theta,\Lambda})$ over $[0,1]$ can be covered by at most $m\bigl(C m^{1-\beta}+1\bigr)^2=m^{3-2\beta}\bigl(C+m^{\beta-1}\bigr)^2$
		cubes of side length \( 1/m \). Therefore,
		\[
		\dim_H G(f_{\Theta,\Lambda})
		\le \overline{\dim}_B G(f_{\Theta,\Lambda})
		\le \limsup_{m\to\infty}
		\frac{\log\bigl(m^{3-2\beta}(C+m^{\beta-1})^2\bigr)}{-\log(1/m)}
		=3-2\beta
		\] where $\overline{\dim}_B$ denotes the upper box dimension, proving the lemma.
	\end{proof}


\bigskip	
 
\begin{theorem}\label{thm2.1}
 	$\dim_H f_{\Theta,\Lambda}([0,1]) = \min\{2, \frac{1}{\beta}\}$ a.s.
 \end{theorem}
 	
 \begin{proof}
 	 Trivially, $\dim_H f_{\Theta,\Lambda}([0,1]) \le 2$, and by the H\"older continuity, we have
 	\[
 	\dim_H f_{\Theta,\Lambda}([0,1]) \le \frac{1}{\beta} \dim_H([0,1]) = \frac{1}{\beta}.
 	\]
 	That proves $\dim_H f_{\Theta,\Lambda}([0,1]) \le \min\{2, \frac{1}{\beta}\}$. For the converse inequality, we let  \( {\mathcal L} \) be the Lebesgue measure on \([0,1]\). Its restriction to $[0, \frac{1}{2b^5}]$  induces the push-forward measure $$\nu := \mathcal{L}|_{[0, \frac{1}{2b^5}]} \circ f_{\Theta,\Lambda}^{-1}$$  which is supported on $f_{\Theta,\Lambda}([0, \frac{1}{2b^5}])$.  Denote by $B=[0, \frac{1}{2b^5}]\times [0, \frac{1}{2b^5}]$.	The Fourier transform of $\nu$ is
 	\[
 	\hat{\nu}(\xi) = \int_{\mathbb{R}^2} e^{-2\pi i \langle \xi, w \rangle} dw,
 	\quad \text{where} \quad  \xi = (\xi_1, \xi_2) \in \mathbb{R}^2.
 	\]
 	Then
 \begin{eqnarray*}
 	 	|\hat{\nu}(\xi)|^2
 	 &=& \int_{B} e^{-2\pi i \langle \xi, f_{\Theta,\Lambda}(x) - f_{\Theta,\Lambda}(y) \rangle} \, d(x,y) \\
 	 &=& \int_{B} e^{-2\pi i (\xi_1 a_1(x,y) + \xi_2 a_2(x,y))} \, d(x,y) \\ 	 
 	 &=& \int_{B} e^{-2\pi i \xi_1 a_1(x,y)} e^{-2\pi i \xi_2 a_2(x,y)} \, d(x,y).
 \end{eqnarray*}
 
 
 	Taking expectation:
 \begin{eqnarray*}
 	\mathbb{E}\left(|\hat{\nu}(\xi)|^2\right)
 	&=& \int_{B} \mathbb{E}\left(e^{-2\pi i \xi_1 a_1(x,y)} e^{-2\pi i \xi_2 a_2(x,y)}\right)\, d(x,y) \\
 	&=& \int_{B} \mathbb{E}\left(e^{-2\pi i \xi_1 a_1(x,y)}\right) \mathbb{E}\left(e^{-2\pi i \xi_2 a_2(x,y)}\right) \, d(x,y).
 \end{eqnarray*}
 	
By substituting the expressions for \(a_1(x,y), a_2(x,y)\), we obtain
 	\[
 	e^{-2\pi i \xi_1 a_1(x,y)}
  	= \prod_{n=0}^\infty e^{2\pi i \xi_1 \varphi_n \sin(\pi(b^n x + b^n y + 2\theta_n))},
 	\]
 	\[
 	e^{-2\pi i \xi_2 a_2(x,y)}
 	= \prod_{n=0}^\infty e^{-2\pi i \xi_2 \varphi_n \cos(\pi(b^n x + b^n y + 2\lambda_n))}.
 	\]
 	Thus,
  \begin{eqnarray*}
 	\mathbb{E}\left(e^{-2\pi i \xi_1 a_1(x,y)}\right)
 	&=& \prod_{n=0}^\infty \mathbb{E}\left(e^{2\pi i \xi_1 \varphi_n \sin(\pi(b^n x + b^n y + 2\theta_n))}\right) \\
    &=& \prod_{n=0}^\infty \int_0^1 e^{2\pi i \xi_1 \varphi_n \sin(2\pi \theta)} d\theta\\
 	&=& \prod_{n=0}^\infty J_0\left(2\pi \xi_1 \varphi_n\right),
 \end{eqnarray*}
 	where $$J_0(x) := \int_0^1 e^{i x \sin(2\pi \theta)} d\theta = \int_0^1 e^{i x \cos(2\pi \theta)} d\theta$$ is the Bessel function satisfying $J_0(-x)=J_0(x)$ (see \cite {Mattila2015} for more details).
 	
 	Similarly,
 	\[
 	\mathbb{E}\left(e^{-2\pi i \xi_2 a_2(x,y)}\right)
 	= \prod_{n=0}^\infty J_0\left(2\pi \xi_2 \varphi_n\right).
 	\]
 	Therefore,
 	\[
 	\mathbb{E}\left(|\hat{\nu}(\xi)|^2\right)
 	= \int_B \prod_{n=0}^\infty \left(J_0(2\pi \xi_1 \varphi_n) J_0(2\pi \xi_2 \varphi_n)\right)\, d(x,y).
 	\]
 	
 For any $0<s<2$, the $s$-energy of $\nu$  is defined by
 $$I_s(\nu)=\int_{\R^2}\int_{\R^2}\frac{d\nu(\xi)d\nu(\eta)}{|\xi-\eta|^s}= \gamma(2,s)\int_{\mathbb{R}^2} |\xi|^{s-2}|\hat{\nu}(\xi)|^2 d\xi$$ where $\gamma(2,s)=\pi^{s-1}\frac{\Gamma(1-s/2)}{\Gamma(s/2)}$ (see details in \cite{Mattila2015}).  By using Fubini-Tonelli theorem and  the inequality $|\xi| \ge |\xi_1|  |\xi_2|$, we can estimate the	expectation of the $s$-energy of $I_s(\nu)$:
 	 \begin{eqnarray}
 	\mathbb{E}(I_s(\nu))
 	&=& \gamma(2,s)\int_{\mathbb{R}^2} |\xi|^{s-2} \mathbb{E}\left(|\hat{\nu}(\xi)|^2\right) d\xi  \nonumber \\
 	&=& \gamma(2,s)\int_{\mathbb{R}^2} |\xi|^{s-2} \int_B \prod_{n=0}^\infty \left(J_0(2\pi \xi_1 \varphi_n) J_0(2\pi \xi_2 \varphi_n)\right) \, d(x,y) d\xi   \nonumber\\
 	&\le& \gamma(2,s)\int_{\mathbb{R}^2} |\xi_1|^{\frac{s-2}{2}} |\xi_2|^{\frac{s-2}{2}} \int_B \prod_{n=0}^\infty \left(J_0(2\pi \xi_1 \varphi_n) J_0(2\pi \xi_2 \varphi_n)\right) \, d(x,y) d\xi    \nonumber  \\
 	&=& \gamma(2,s)\int_B \left(\int_{\mathbb{R}} |t|^{\frac{s-2}{2}} \prod_{n=0}^\infty (J_0(2\pi t \varphi_n))dt\right)^2 \, d(x,y)   \nonumber  \\
 	&=& \gamma(2,s)\sum_{k=5}^{\infty}\int_{B_k} \left(\int_{\mathbb{R}} |t|^{\frac{s-2}{2}} \prod_{n=0}^\infty (J_0(2\pi t \varphi_n))dt\right)^2 \, d(x,y)  \label{eq-expectation}
  \end{eqnarray}
  where $$B_k=\left\{(x,y)\in [0,1]^2: \frac{1}{2b^{k+1}}<|x-y|\le \frac{1}{2b^k}\right\} \quad \text{and}\quad B=\bigcup_{k=5}^\infty B_k.$$    
  It is easy to verify that the Lebesgue measure of $B_k$ is bounded by $\mathcal{L}^2(B_k)\le C_0 b^{-k}$ for some $C_0>0$.    For any $k\ge 5$ and $(x,y)\in B_k$, when $n=k-4,\dots, k$,  we know that
   $$\pi b^n|x-y|\in \left(\frac{\pi}{2b^{k+1-n}},\frac{\pi}{2b^{k-n}}\right]\subset \left(0,\frac{\pi}{2}\right].$$  
   By using the global esitmate of Bessel function  $|J_0(x)|\le (1+|x|^2)^{-1/4}$ (refer to \cite {Mattila2015}), we have 
  \begin{eqnarray*} 
  |J_0(2\pi t \varphi_n)| &\le& |2\pi t \varphi_n|^{-\frac12} \\
  &=& (4\pi t)^{-\frac12}b^{\frac{\beta n}{2}}|\sin(\pi b^n(x-y))|^{-\frac12} \\
  &\le& (2\pi^2t)^{-\frac12}b^{\frac{k+1+(\beta-1)n}{2}}.
  \end{eqnarray*}
  It follows that 
  \begin{eqnarray*} 
  	\prod_{n=k-4}^k |J_0(2\pi t \varphi_n)| &\le& (2\pi^2t)^{-\frac52}b^{\frac{4k\beta-10\beta+15}{2}}.
   \end{eqnarray*}
 	
 	 Denote by
 	 \begin{eqnarray*}
 	 	\Phi(x,y) &:=& \int_{\mathbb{R}} |t|^{\frac{s-2}{2}} \prod_{n=0}^\infty J_0(2\pi t \varphi_n)dt \\
 	 	&=& \left(\int_{|t|\le b^{\beta k}}+\int_{|t|> b^{\beta k}}\right) |t|^{\frac{s-2}{2}} \prod_{n=0}^\infty J_0(2\pi t \varphi_n)dt  \\
 	 	&:=& I_1+I_2.
 	 \end{eqnarray*}
 	After some computations, we get
  $$|I_1|\le\int_{|t|\le b^{\beta k}}|t|^{\frac{s-2}{2}} \prod_{n=0}^\infty |J_0(2\pi t \varphi_n)|dt\le 2\int_{0}^{b^{\beta k}}t^{\frac{s-2}{2}}dt=\frac{4}{s} b^{\frac{1}{2}\beta sk},$$
 	\begin{eqnarray*} 
 	|I_2| &\le& \int_{|t|> b^{\beta k}}|t|^{\frac{s-2}{2}} \prod_{n=0}^\infty |J_0(2\pi t \varphi_n)|dt \\
 	&=& 2\int_{b^{\beta k}}^\infty t^{\frac{s-2}{2}} \prod_{n=0}^\infty |J_0(2\pi t \varphi_n)|dt \\
 	&\le& 2\int_{b^{\beta k}}^\infty t^{\frac{s-2}{2}} \prod_{n=k-4}^k |J_0(2\pi t \varphi_n)|dt \\
 	&\le& 2\int_{b^{\beta k}}^\infty t^{\frac{s-2}{2}}(2\pi^2t)^{-5/2}b^{\frac{4k\beta-10\beta+15}{2}} dt \\
 	&=& 2(2\pi^2)^{-5/2}b^{\frac{1}{2}(5k\beta-10\beta+15)}\int_{b^{\beta k}}^\infty t^{\frac{s-7}{2}} dt  \\
 	&\le& C_1 b^{\frac{1}{2} \beta ks}
 	  \end{eqnarray*}
 where $C_1>0$ is a constant.  Hence 
 \begin{equation}\label{eq-Phi}
 	 |\Phi(x,y)|\le |I_1|+|I_2|\le C_2 b^{\frac{1}{2} \beta ks}, \quad\text{where}\quad C_2 :=C_1+\frac 4s.
 \end{equation} 
 
 From \eqref{eq-expectation} \eqref{eq-Phi}  it follow that
  \begin{eqnarray*}
 	\mathbb{E}(I_s(\nu))&\le& \gamma(2,s) \sum_{k=5}^{\infty}\int_{B_k}(\Phi(x,y))^2 \, d(x,y) \\
 	&\le&  \gamma(2,s)\sum_{k=5}^{\infty} (C_2)^2 b^{\beta ks} \mathcal{L}^2(B_k) \\
 	&\le&  \gamma(2,s)C_3\sum_{k=5}^{\infty} b^{(\beta s-1)k}.
 \end{eqnarray*}
 
 If  $\beta\in (0, \frac12]$, then for any $0<s<2$, we have $\beta s<1$, and $\mathbb{E}(I_s(\nu))<\infty$, which implies that $\dim_H f_{\Theta,\Lambda}([0,\frac{1}{2b^5}])\ge 2$ a.s.
 
 If $\beta\in (\frac12, 1)$, then for any $0<s<\frac {1}{\beta} \  (<2)$, we have  $\beta s<1$, and $\mathbb{E}(I_s(\nu))<\infty$, which implies that $\dim_H f_{\Theta,\Lambda}([0,\frac{1}{2b^5}])\ge \frac{1}{\beta}$ a.s. Therefore, it concludes that  $\dim_H f_{\Theta,\Lambda}([0,1])\ge \dim_H f_{\Theta,\Lambda}([0,\frac{1}{2b^5}])\ge \min\{2, \frac{1}{\beta}\}$ a.s.
 \end{proof}
 
  
\bigskip
  
 \begin{theorem} 
 	The Hausdorff dimension of the graph $f_{\Theta,\Lambda}$ satisfies
 	$$\dim_H G(f_{\Theta,\Lambda}) =\dim_B G(f_{\Theta,\Lambda}) = \min\left\{\frac{1}{\beta}, \, 3 -2\beta\right\} \quad \text{a.s.}$$
 \end{theorem}
 
	
\begin{proof} 
	Define the projection  $P: G(f_{\Theta,\Lambda}) \to f_{\Theta,\Lambda}([0,1])$ from the graph onto the image by 
	$$P(x, f_{\Theta,\Lambda}(x))= f_{\Theta,\Lambda}(x), \quad x\in [0,1].$$
	Trivially $P$ is a Lipschitz function.  Hence $\dim_H G(f_{\Theta,\Lambda})\ge \dim_H f_{\Theta,\Lambda}([0,1]).$
	
    If $\beta\in [\frac12, 1]$, 	it follows from  Lemma \ref{lem-lowerbound} and Theorem \ref{thm2.1} that
	$$ \frac {1}{\beta} \ge \dim_H G(f_{\Theta,\Lambda})\ge \dim_H f_{\Theta,\Lambda}([0,1])=\frac{1}{\beta} \quad \text{a.s.}$$
	That proves  $\dim_H G(f_{\Theta,\Lambda}) = \frac {1}{\beta}$ a.s.
	
	 If $\beta\in (0, \frac12)$,  we let   \( \mu:= \mu_{\Theta,\Lambda} \) be the induced measure on $\mathbb{R}^3$ through the Lebesgue measure  \( {\mathcal L} \) and $f_{\Theta,\Lambda}$ in the following way: for any Borel set \( V \subset \mathbb{R}^3 \), define
	\begin{equation*}
		\mu (V) = {\mathcal L}\bigl(\{t \in [0,1] : (t, f_{\Theta,\Lambda}(t)) \in V\}\bigr).
	\end{equation*} 
	The measure $\mu$ is supported on the graph  $G(f_{\Theta,\Lambda})$.  
	
	For $s>0$, the \(s\)-energy of \( \mu \)  is given by
	\begin{eqnarray*}
			I_s(\mu)
			&=& \int_{\mathbb{R}^3} \int_{\mathbb{R}^3}\frac{d\mu(w)d\mu(z)}{|w-z|^s} \\
			&=& \int_{[0,1]^2} \frac{d(x,y)}{\bigl((x-y)^2 + |f_{\Theta,\Lambda}(x)-f_{\Theta,\Lambda}(y)|^2\bigr)^{s/2}}.
	\end{eqnarray*}
	
	We will prove that for any \( s\in (2, 3-2\beta) \), the \(s\)-energy \( I_s(\mu) \) is finite for almost every pair of phase sequences \( (\Theta,\Lambda) \in H \times H \). This implies that the Hausdorff dimension of the graph of \( f_{\Theta,\Lambda} \) is at least \( s \). By letting \( s \to 3-2\beta \), we conclude that
	\[
	\dim_H G(f_{\Theta,\Lambda}) \ge 3-2\beta
	\]
	holds for almost every \( (\Theta,\Lambda) \). Hence the desired result follows by Lemma \ref{lem-lowerbound}.
	
	In order to show $I_s(\mu)<\infty $ a.s., it suffices to prove that the expectation \( \mathbb{E}(I_s(\mu)) < \infty \). By Fubini-Tonelli theorem, we have
	\begin{equation*}
		\mathbb{E}(I_s(\mu))= \int_{[0,1]^2} \mathbb{E}\left(\frac{1}{\bigl((x-y)^2 + |f_{\Theta,\Lambda}(x)-f_{\Theta,\Lambda}(y)|^2\bigr)^{s/2}}\right) \, d(x,y).
	\end{equation*}
	
	We claim that there exists $c_0>0$ such that for \( 0 < |x-y| \le \frac{1}{2b^2} \),  
	$$\mathbb{E}\left(\frac{1}{\bigl((x-y)^2 + |f_{\Theta,\Lambda}(x)-f_{\Theta,\Lambda}(y)|^2\bigr)^{s/2}}\right)\le c_0|x-y|^{2-s-2\beta}.$$
	Since \( 2 < s < 3-2\beta \), it concludes that $\mathbb{E}(I_s(\mu))<\infty$.
	
	Fix \( x,y \in [0,1] \) satisfying \( 0 < |x-y| \le \frac{1}{2b^2}\), and define
	$$X = |f_{\Theta,\Lambda}(x)-f_{\Theta,\Lambda}(y)|^2.$$
	Regarding \( X \) as a random variable depending on \( (\Theta,\Lambda) \), it can be verified that \( X \) has a bounded probability density function \( h \). Then
	\begin{eqnarray*}
			&& \mathbb{E}\left(\frac{1}{\bigl((x-y)^2 + |f_{\Theta,\Lambda}(x)-f_{\Theta,\Lambda}(y)|^2\bigr)^{s/2}}\right) \\
			&= &\int_0^\infty \frac{h(t)\,dt}{\bigl((x-y)^2 + t\bigr)^{s/2}} \\
			&= &\int_0^\infty \frac{h\bigl(|x-y|^2 \tau\bigr) |x-y|^2 d\tau}{|x-y|^s (1+\tau)^{s/2}} \\
			&\le &\Bigl(\sup_{t\ge 0} h(t)\Bigr) \, |x-y|^{2-s} \int_0^\infty \frac{d\tau}{(1+\tau)^{s/2}}.
	\end{eqnarray*}
	 Hence to prove the claim, we need to show that there exists a constant \( c_1>0 \), independent of \( x,y \), such that
	\begin{equation}\label{eq-claim}
	h(t) \le c_1|x-y|^{-2\beta}.
	\end{equation}
	
	By using the notation in \eqref{eq-phi-n} and \eqref{eq- vector-valued funct.}, $X$ can be written as
	$$X = (a_1(x,y))^2 +(a_2(x,y))^2 := X_1^2 + X_2^2,$$
	where    
	   $$X_1 = \sum_{n=0}^\infty  \varphi_n \sin\left(\pi\left(b^n x + b^n y +2\theta_n\right)\right)$$ 
	and 
	   $$X_2 = \sum_{n=0}^\infty  \varphi_n \cos\left(\pi\left(b^n x + b^n y +2\lambda_n\right)\right).$$
   Since \( \theta_n \sim \mathrm{Uniform}[0,1] \) are i.i.d., each \( \varphi_n \sin\left(\pi\left(b^n x + b^n y +2\theta_n\right)\right) \) has density function
	\begin{equation*}
		h_n(z) =
		\begin{cases}
			\dfrac{1}{\pi\sqrt{{\varphi_n}^2-z^2}} & |z| < |\varphi_n|, \\
			0 & |z| \ge |\varphi_n|.
		\end{cases}
	\end{equation*}
 The density function of $X_1$ is the infinite convolution 
 $$h_{X_1} = h_0 * h_1 * h_2 * \cdots.$$ 
 Since the supremum of a probability density is nonincreasing under convolution, it follows that	
	\[
	\sup_t h_{X_1}(t) \le \sup_t \bigl(h_{k-2} * h_{k-1} * h_k\bigr)(t)
	\]
	for any fixed \( k \ge 2 \). 
	
	Choose an integer \( k \) such that $\frac{1}{2b^{k+1}} < |x-y| < \frac{1}{2b^k}$. 	Then
	\begin{equation*}
		\frac{\pi}{2b^3} < \Bigl|\pi b^{k-2}(y-x)\Bigr|\le \Bigl|\pi b^k(y-x)\Bigr| \le \frac{\pi}{2}.
	\end{equation*}
	So for \( n = k-2,k-1,k \),
	\begin{equation*}
		|\varphi_n|=|2b^{-\beta n}\sin(\pi b^n(x-y))| > 2\sin\Bigl(\frac{\pi}{2b^3}\Bigr) b^{-k\beta}
		> 2^{1+\beta}\sin\Bigl(\frac{\pi}{2b^3}\Bigr) |x-y|^\beta.
	\end{equation*}
	
	Let \( \|\cdot\|_p \) denote the \( L^p \)-norm. For \( p=3/2 \),
	\begin{eqnarray*}
		\|h_n\|_{\frac32}
			&=& \left( \int_{-|\varphi_n|}^{|\varphi_n|} \frac{dz}{\bigl(\pi\sqrt{{\varphi_n}^2-z^2}\bigr)^{\frac32}} \right)^{\frac23} \\
			&=& \pi^{-1} \left( \int_{-|\varphi_n|}^{|\varphi_n|} ({\varphi_n}^2-z^2)^{-\frac34}dz \right)^{\frac23} \\
			&=& \pi^{-1} I^{\frac23} |\varphi_n|^{-\frac13},
	\end{eqnarray*}
	where \( I = \int_{-1}^1 (1-u^2)^{-\frac34}du < \infty \) is a universal constant. Thus for \( n=k-2,k-1,k \),
	\[
	\|h_n\|_{\frac32} \le c_2 |x-y|^{-\frac{\beta}{3}}
	\]
	where \( c_2 \) depends only on \( b \). 	By Young's inequality and H\"older's inequality,
	\begin{eqnarray*}
		\bigl(h_{k-2} * h_{k-1} * h_k\bigr)(t)
			&\le& \|h_{k-2}\|_{\frac32} \|h_{k-1} * h_k\|_3 \\
			&\le&  \|h_{k-2}\|_{\frac32} \|h_{k-1}\|_{\frac32}\|h_k\|_{\frac32} \\
			& \le& c_3 |x-y|^{-\beta}
	\end{eqnarray*} 
	where $c_3 :=(c_2)^3$. Hence $h_{X_i}(t) \le c_3 |x-y|^{-\beta}$. The same bound holds for  $h_{X_2}$, the density function of $X_2$.
	
	Finally, we compute the density function of $X= X_1^2 + X_2^2$. For \( i=1,2 \), the density of \( X_i^2 \) is
	\begin{equation*}
		h_{X_i^2}(t) =
		\begin{cases}
			\dfrac{h_{X_i}(\sqrt{t})+h_{X_i}(-\sqrt{t})}{2\sqrt{t}}, & t\ge 0, \\
			0, & t<0.
		\end{cases}
	\end{equation*}
	Since \( X_1 \) and \( X_2 \) are independent, the density function of \( X \) satisfies
	\begin{eqnarray*}
			h(t) &=& \bigl(h_{X_1^2} * h_{X_2^2}\bigr)(t) \\
			&=& \int_0^t \frac{h_{X_1}(\sqrt{t-s})+h_{X_1}(-\sqrt{t-s})}{2\sqrt{t-s}}
			\cdot \frac{h_{X_2}(\sqrt{s})+h_{X_2}(-\sqrt{s})}{2\sqrt{s}} ds \\
			&\le& ({c_3})^2|x-y|^{-2\beta} \int_0^t \frac{ds}{\sqrt{s(t-s)}} \\
			&=& ({c_3})^2\pi |x-y|^{-2\beta}.
	\end{eqnarray*}
By letting $c_1 :=\pi ({c_3})^2$, we obtain the inequality \eqref{eq-claim}, hence prove the theorem.
\end{proof}


	
\bigskip

\noindent {\bf Acknowledgements:} 
The first draft of this paper was completed while the first author was visiting the Department of Mathematics at The Chinese University of Hong Kong in the summer of 2023. He would like to express his sincere gratitude to Prof. De-Jun Feng for his hospitality. The authors are also deeply grateful to Dr. Peng-Fei Zhang for many valuable discussions and suggestions.
	
\bigskip

\begin{thebibliography}{99}
	
	\bibitem{Baranski2002}
	Barański K. On the complexification of the Weierstrass non-differentiable function.
	\textit{Annales Academiae Scientiarum Fennicae Mathematica}, 2002, 27(2): 325--339.
	
	\bibitem{Baranski2012}
	Barański K. On the dimension of graphs of Weierstrass-type functions with rapidly growing frequencies.
	\textit{Nonlinearity}, 2012, 25(1): 193--209.
	
	\bibitem{Baranski2014}
	Barański K, Bárány B, Romanowska J. On the dimension of the graph of the classical Weierstrass function.
	\textit{Advances in Mathematics}, 2014, 265: 32--59.
	
		
	\bibitem{Falconer2014}
	Falconer K J. \textit{Fractal Geometry: Mathematical Foundations and Applications}.
	3rd ed. Chichester: John Wiley \& Sons, 2014.
	
	\bibitem{Hardy1916}
	Hardy G H. Weierstrass's non-differentiable function.
	\textit{Transactions of the American Mathematical Society}, 1916, 17(3): 301--325.
	
	\bibitem{Heurteaux2003}
	Heurteaux Y. Dimension of graphs of random Weierstrass functions.
	\textit{Annales de l'Institut Henri Poincaré, Probabilités et Statistiques}, 2003, 39(2): 273--296.
	
	\bibitem{Hunt1998}
	Hunt B R. The Hausdorff dimension of graphs of Weierstrass functions.
	\textit{Proceedings of the American Mathematical Society}, 1998, 126(3): 791--800.
	
	\bibitem{Kahane1963}
	Kahane J-P, Weiss M, Weiss G. On lacunary power series.
	\textit{Arkiv för Matematik}, 1963, 5: 1--26.
	
	\bibitem{Kaplan1984}
	Kaplan J L, Mallet-Paret J, Yorke J A. The Lyapunov dimension of a strange attractor.
	\textit{Physica D: Nonlinear Phenomena}, 1984, 13(3): 457--476.
	
	\bibitem{LaiLauZhang2026}
	Lai C-K, Lau K-S, Zhang P-F. Hausdorff dimension of images and graphs of some random complex series. arXiv:2603.05986 [math.CA], 2026.
	
	\bibitem{Ledrappier1992}
	Ledrappier F. On the dimension of the graph of the Weierstrass function.
	\textit{Studia Mathematica}, 1992, 101(3): 233--246.
	
	\bibitem{Mandelbrot1977}
	Mandelbrot B B. \textit{Fractals: Form, Chance, and Dimension}.
	San Francisco: W. H. Freeman, 1977.
	
	\bibitem{Mauldin1986}
	Mauldin R D, Williams S C. On the Hausdorff dimension of some graphs.
	\textit{Transactions of the American Mathematical Society}, 1986, 298(2): 793--803.
	
	\bibitem{Mattila2015}
	Mattila P. \textit{Fourier Analysis and Hausdorff Dimension}.
	Cambridge Studies in Advanced Mathematics 150. Cambridge: Cambridge University Press, 2015, pp. 1--440.
	
	\bibitem{Przytycki1989}
	Przytycki F, Urbański M. On the Hausdorff dimension of some fractal sets.
	\textit{Studia Mathematica}, 1989, 93(2): 155--186.
	
	\bibitem{Ren2021}
	Ren H, Shen W. A dichotomy for the Weierstrass-type functions.
	\textit{Inventiones Mathematicae}, 2021, 225(2): 549--602.
	
	\bibitem{Ren2023}
	Ren H. Box dimension of the graphs of the generalized Weierstrass-type functions.
	\textit{Discrete and Continuous Dynamical Systems -- Series A}, 2023, 43(10): 3830--3838.
	
	\bibitem{RenShen2024}
	Ren H, Shen W. The high-dimensional Weierstrass functions. arXiv:2404.06778 [math.CA], 2024.
	
	\bibitem{Rom2014}
	Romanowska J. Measure and Hausdorff dimension of randomized Weierstrass-type functions. \textit{Nonlinearity}, 2014, 27: 787--801.
	
	\bibitem{Salem1945}
	Salem R, Zygmund A. Lacunary power series and Peano curves.
	\textit{Duke Mathematical Journal}, 1945, 12: 569--578.
	
	\bibitem{Shen2018}
	Shen W. Hausdorff dimension of the graphs of the classical Weierstrass functions.
	\textit{Mathematische Zeitschrift}, 2018, 289(1--2): 223--266.
	
	\bibitem{Solomyak1995}
	Solomyak B. On the random series $\sum \pm\lambda^n$ (an Erdős problem).
	\textit{Annals of Mathematics}, 1995, 142(3): 611--625.
	
	\bibitem{Weierstrass1872}
	Weierstrass K. Über continuirliche Functionen eines reellen Arguments, die für keinen Werth des letzteren einen bestimmten Differentialquotienten besitzen.
	\textit{Mathematische Werke}, Vol. 2, 1872, pp. 71--74.
	
\end{thebibliography}
	
\end{document}
